\newcommand{\tipoRequerimento}{Solicitação de Histórico de Demanda da IP Estimada}

\section{FUNDAMENTAÇÃO LEGAL DA REQUISIÇÃO}

\begin{enumerate}[resume]

\item Do embasamento legal:
\Citacao{Art. 399. O consumidor e demais usuários podem requerer informações, solicitar e cancelar serviços, encaminhar reclamações, elogios, sugestões, denúncias e solicitar o encerramento contratual nos canais de atendimento disponibilizados pela distribuidora. (Redação dada pela REN ANEEL 1.042, de 20.09.2022) § 1º O consumidor e demais usuários podem requerer informações, esclarecer dúvidas, encaminhar sugestões, reclamações, denúncias, elogios e sugestões à Ouvidoria da distribuidora, se houver, à agência estadual conveniada ou, na inexistência desta, à ANEEL, observado o art. 421 e o art. 423. § 2º O atendimento deve ser classificado e registrado de acordo com opção do consumidor e demais usuários quanto à natureza de sua demanda, entre informação, reclamação, serviço, denúncia, elogio ou sugestão, conforme os procedimentos e tipologia estabelecidos pela ANEEL. Art. 400. Independentemente do canal escolhido pelo consumidor e demais usuários, recomenda-se que a distribuidora solucione as demandas no primeiro contato, implementando mecanismos de medição e de acompanhamento da eficiência operacional e da satisfação do atendimento. Art. 401. Na recepção da demanda do consumidor e demais usuários, a distribuidora deve considerar como reclamação: I - demanda que expresse insatisfação do consumidor e demais usuários em relação aos serviços e produtos recebidos ou em relação à atuação da distribuidora; II - demanda que tenha por objetivo o reconhecimento de um direito, ou a correção de ato que cause lesão ou ameaça de lesão a direitos do consumidor e demais usuários; III - comunicação relacionada à interrupção do fornecimento de energia elétrica; IV - solicitação de ressarcimento de danos elétricos em equipamentos V - solicitação que implique correção de faturamento, ainda que realizada durante o atendimento; e VI - todas aquelas provenientes de plataforma “Consumidor.gov.br”, mantida pelo Ministério da Justiça. Art. 406. A distribuidora deve responder a demanda, preferencialmente, pelo canal utilizado para o protocolo, ou por outro canal previamente escolhido pelo consumidor e demais usuários para o relacionamento com a distribuidora, observado o art. 416. § 1º Quando der a resposta por meio de canal telefônico, a distribuidora deve esclarecer ao consumidor e demais usuários sobre o seu direito de optar pelo recebimento das informações por escrito, por meio de correspondência ou por meio eletrônico. § 2º A distribuidora pode disponibilizar as informações no espaço reservado na internet ou em canal alternativo, desde que o próprio consumidor ou demais usuários tenham optado por essa forma de relacionamento. § 3º A distribuidora deve responder de forma clara, objetiva, conclusiva e abordar todos os pontos da demanda do consumidor e demais usuários, utilizando linguagem simples, comum, de fácil compreensão e evitando o uso de termos técnicos, explicando-os quando necessário. Art. 407. A distribuidora deve prestar as informações solicitadas pelo consumidor e demais usuários de forma imediata.}

\item Cumpre registrar que (i) ao negar ou indeferir o pedido sem fundamento, (ii) ao responder de forma parcial sem observar todas as informações solicitadas, (iii) ao apresentar resposta divergente do que consta no sistema informático da Enel, a ENEL pode estar cometendo as seguintes irregularidades, passíveis de sanção por parte da ANEEL:

\Citacao{Art. 9º Constitui infração do Grupo I: III - deixar de prestar informações aos consumidores ou usuários, quando solicitado ou conforme determinado nas disposições legais, regulamentares ou contratuais; VII - deixar de registrar ou de analisar as ocorrências nos seus sistemas de distribuição, transmissão ou geração; [...] Art. 10. Constitui infração do Grupo II: I - deixar de manter registro atualizado das reclamações e solicitações dos consumidores; [...] Art. 11. Constitui infração do Grupo III: [...] VII - deixar de cumprir ao disposto nos Procedimentos de Distribuição; VIII - deixar de cumprir ao disposto nos Procedimentos de Rede; IX - deixar de cumprir ao disposto nos Procedimentos de Regulação Tarifária; X - deixar de cumprir ao disposto nas Condições Gerais de Fornecimento de Energia Elétrica; [...] XIV - deixar de cumprir ao disposto em resoluções da ANEEL; [...] XXII - enviar à ANEEL documento ou informação com conteúdo incorreto, resultando em prejuízo às decisões da Agência ou à divulgação de informações. [...] Art. 12. Constitui infração do Grupo IV: I - descumprir às disposições legais, regulamentares e contratuais relativas: III - deixar de atender pedido de serviços nos prazos e nas condições estabelecidas na legislação ou no contrato; [...] XIII - impor ônus para o solicitante ou consumidor na prestação do serviço público de distribuição de energia elétrica em desacordo com as disposições legais ou regulamentar [...] Art. 13. Constitui infração do Grupo V: [...] XVII - fornecer documentos ou informações falsas à ANEEL;}

\item Por fim, cumpre esclarecer que ao fornecer informação diferente da que consta em seu sistema informático, pode configurar infração mais grave. 

\item Portanto, a solicitação tem previsão legal, conforme explicitado. Desta forma, não cabe respostas de indeferimentos vazias, sem embasamento legal, ou resposta de deferimento parcial com envio de parte das informações solicitadas.
\end{enumerate}

\section{DOS REQUERIMENTOS}

\begin{enumerate}[resume]

\item Solicita-se que a conssecionária apresente, o histórico e conteúdo das demandas (reclamações/solicitações) dos últimos \textbf{114 meses} relacionadas à Unidade Consumidora (UC) em epígrafe, contendo, no mínimo, as seguintes informações:  
  \begin{enumerate}
    \item Número do protocolo.  
    \item Classificação da solicitação/reclamação.  
    \item Avaliação de procedência ou improcedência realizada pela distribuidora.  
    \item Datas de solicitação e de solução da distribuidora, tempo total transcorrido prazo regulamentar.  
    \item Providências adotadas pela distribuidora.  
    \item Conteúdo da resposta e providências adotadas pela distribuidora.  
    \item Demais informações relacionadas às demandas. 
  \end{enumerate}
\item Estão incluídas no rol de demandas acima, dentre outras, as seguintes:
  \begin{enumerate}
    \item Solicitação de realização de alteração cadastral,
    \item Fornecimento de orçamento estimado,
    \item Informação de resultado de análise de projeto,
    \item Informação de resultado de classificação tarifária,
    \item Respostas ao cliente de reclamações em geral,
    \item Verificação ou retirada do equipamento danificado,
    \item Entrega dos contratos e demais documentos,
    \item Conclusão das obras de conexão,
    \item Vistoria e instalação de medição,
    \item Alteração de titularidade,
    \item Vistorias realizadas,
    \item Substituição de medidor realizada e prazos de realização,
    \item Reclamação de variação de consumo,
    \item Solicitação de acréscimo ou decréscimo de carga,
    \item Reclamação de conta (fatura) não entregue,
    \item Solicitação e resposta de análise de viabilidade técnica,
    \item Solicitação de equipamentos de medição.
  \end{enumerate}
\item Solicita-se que a referida solicitação seja cumprida integralmente dentro do prazo legal estabelecido de 3 (três) dias úteis, conforme estabelece o parágrafo único do art. 418 da REN 1.000/2021.
\Citacao{Art. 418. O consumidor e demais usuários têm direito ao conteúdo do histórico de suas demandas dos últimos 10 anos, observado o art. 670. Parágrafo único. A distribuidora deve informar em até 3 dias úteis, contados a partir da solicitação, no mínimo, as seguintes informações: I - número do protocolo; II - classificação da demanda, conforme tipologia definida pela ANEEL; III - avaliação de procedência ou improcedência realizada pela distribuidora, caso aplicável; IV - datas de solicitação e de solução da distribuidora, tempo total transcorrido e prazo regulamentar; V - conteúdo da resposta e providências adotadas pela distribuidora; (Redação dada pela REN ANEEL 1.042, de 20.09.2022) VI - valores de compensação creditados na fatura pelo descumprimento do prazo regulamentar e mês do crédito; VII - demais informações relacionadas à demanda. (Redação dada pela REN ANEEL 1.042, de 20.09.2022)}

\end{enumerate}
